%!TEX program = lualatex

\documentclass[portuguese, 11pt, a4paper, titlepage, oneside]{article}

\usepackage{geometry}
\usepackage[T1]{fontenc}
\usepackage{graphicx}
\usepackage{parskip}
\usepackage{amssymb}
\usepackage[portuguese]{babel}
\usepackage[framed, numbered]{matlab-prettifier}
\usepackage{subcaption}
\usepackage{wrapfig}
\usepackage{hyperref}

% Title page config
\title{
  \includegraphics[width=10cm]{assets/deti_horizontal_preto.pdf}
  \vskip 2em
  \huge Aspetos Profissionais e Sociais\\
  da Engenharia Informática
  \vskip 1em
	{\bfseries Trabalho individual 1}
}

\author{Rúben Garrido\\NMec: 107927}
\date{\today}
% Begin document
\begin{document}

% Title
\maketitle

\section{Utilização do ChatGPT pelos alunos e professores, no contexto de APSEI}
\subsection{Introdução}
A inteligência artificial baseada em \textit{Large Language Models} (neste caso, o ChatGPT) revolucionou a forma como os humanos pensam e interagem com a tecnologia. Todas as áreas foram influenciadas, embora a educação tenha tido um dos maiores destaques.

Esta mudança de paradigma deve-se sobretudo à utilização do ChatGPT (bem como de outros LLMs) pelos alunos, por forma a "facilitar" a curva de pensamento necessária para realizar um determinado trabalho. No entanto, o uso excessivo e irresponsável da inteligência artificial levou ao surgimento de respostas meramente criadas por esta, com imprecisões e padrões repetidos.

Aspetos Profissionais e Sociais da Engenharia Informática é uma Unidade Curricular que, ao contrário das restantes UCs do curso, se foca em questões éticas e de posicionamento dos engenheiros informáticos perante a sociedade. Esta visa promover a formação de profissionais conscientes do mundo ao seu redor a nível legislativo, financeiro e social, bem como das implicações da tecnologia no panorama atual.

O objetivo deste trabalho foca-se na perceção da usabilidade do ChatGPT no âmbito do processo de ensino e aprendizagem da unidade curricular APSEI, por forma a um melhor entendimento geral no que concerne ao processo de colaboração ou comprometimento dos processos de ensino.

\subsection{Ética e moral}
À primeira instância, o que me ocorre de imediato é "Não! Seria imoral, enganador...". Dependendo em que situação ou de que forma seria utilizado o ChatGPT, como tudo na sociedade existem pontos positivos e pontos negativos, e certamente que o uso desta ferramenta, na UC de APSEI poderia levantar questões éticas e imorais, uma vez que certos trabalhos apelam à criatividade de cada aluno, que, ao utilizarem a inteligência artificial, poderiam colocar em causa a hipótese ou não de um plágio, estando o aluno a pôr em questão a sua ética profissional.

Se, neste momento, para realizar este trabalho eu recorresse à IA por completo, estaria a reduzir a minha capacidade criativa e original de pensamento, e, moralmente, estaria a enganar o docente, o que eticamente, a nível de integridade académica não é correto. APSEI requer muita criatividade artística de pensamento "fora da caixa", o que, para alguns alunos, poderá ser a "pedra no sapato" do curso. Recorrendo ao ChatGPT, os alunos poderão ver ali o escape e, ao mesmo tempo, a solução para o problema que a maioria enfrenta da escrita. Compete aos alunos terem o discernimento de usarem esta ferramenta de forma útil mas responsável, quanto aos aspetos éticos da unidade curricular.

A integração da inteligência artificial vem trazer algumas ideias novas, em situações onde nem sempre se consegue discernir por nós próprios. Esta permite simular alguns dilemas éticos, e inclusive olhar para um certo problema com uma visão diferente. [\textit{Dizer que estou a utilizar o Gemini para pensar em tópicos}]

Para além disso, a utilização do ChatGPT pode ser benéfica \textbf{no que toca à correção linguística, gramatical e ortográfica}. Uma vez que APSEI é a única unidade curricular do curso que onde é necessário realizar trabalhos de escrita, é comum os alunos não se sentirem tão confortáveis a escrever artigos de opinião, devido à falta de prática.

Assim, e por um lado, o ChatGPT permite corrigir erros gramaticais e reescrever frases menos bem construídas. Por outro, quando utilizada com espírito crítico, pode ser um incentivo aos alunos de modo a que estes desenvolvam \textit{skills} linguísticas mais aprofundadas.

\subsubsection{Aspetos negativos}
Texto

\subsection{Professores}
Texto

\section{Should the teacher tape the classes and make the videos available for the students later?}
\begin{itemize}
  \item If yes, any limitations on taping or on the distributions?
  \item If no, why not?
  \item In both cases, what are the cons of your choice, both from the teacher and students' point of view?
  \item Deve ter microfones a apanhar o áudio de toda a gente?
  \item Deve-se gravar a intervenção dos alunos?
  \item Deve-se partilhar para quem não vai às aulas? Tratar por igual?
  \item Aplica-se a todas as aulas?
  \item O número de alunos influencia a opinião?
  \item Asneiras nas aulas
  \item Contexto geográfico?
\end{itemize}

\section{O que falta abordar}
\begin{itemize}
  \item If yes, should students be expected to use ChatGPT?
  \item If yes, should students be incentivized to use ChatGPT in which way?
  \item If no, why not?
  \item And can/should professors follow the same rules as the students?
\end{itemize}
\end{document}